\scanpage{124}
% archive foliation: 123
(auquel cas les pts correspondants des
$A_X$ \uline{ne sont pas} des sommets topologiques
\unsure{\ldots}) \deleted{\unsure{on}} \deleted{\unsure{en}} \inserted{\unsure{se permet de}}
\unsure{le considère|les considère} comme un ``milieu
d'arête'' des $A_X$.

À l'exclusion de ce cas, on peut dire
que $A$, avec l'ens de ses sommets
choisis (qui \unsure{c.} \deleted{\unsure{A}} \uline{tous} \uline{les} sommets
topologiques, et les pts $0,1,\infty$ \inserted{qui sont $\in A$}, (qu'ils
\unsure{soient ou non} sommets top., et éventuellement
$N$ autres, \unsure{enfin} suivant l'opération
qu'on veut obtenir sur les sommets
qu'on entend donner aux $A_X$), et
le choix des centres des faces fait
en conformité avec la condition $1^\circ$,
et enfin un choix \uline{arbitraire} des
milieux d'arête de $A$ (tous distincts
\unsure{même} de $0,1,\infty$) est \unsure{déduit} d'une
application $\varphi:S\to S$ ramifiée aux
seuls pts $0,1,\infty$, avec ramification
$2$ au dessus de $1$, en prenant l'image
inverse de $[0,1]$ -- et l'image inverse
\deleted{dans} $A_X$ dans $X$ par $f:X\to S$ de
$A$ n'est que $(\varphi f)^{-1}([0,1])$,
\unsure{où} $\varphi f$ est \deleted{ramifiée} un rev. \unsure{en dehors de}
$0,1,\infty$ (car $\varphi^{-1}\{0,1,\infty\}\supset\{0,1,\infty\}$ et
